% Teil II: der durchgerechnete Weg.

% Gedankenschritt: Ränder-Probe (erster und letzter Term) beim Umschreiben.

\subsection*{Gegeben und gesucht}
Umrechnen zwischen ausgeschriebener Summe und Summenzeichen, ohne Werte zu
berechnen.

\subsection*{Ansatz}
\textbf{Das Muster lautet: Wir übersetzen, wir rechnen nicht.} Für jedes
Teilrätsel identifizieren wir (1) den allgemeinen Summanden $a_i$,
(2) den Startwert, bei dem er den ersten Term liefert, und (3) den
Endwert, bei dem er den letzten Term liefert. Die Ränder-Probe am Anfang
und Ende ist der entscheidende Kontrollschritt.

\subsection*{a) $2+3+4+\dots+12$}
Der Summand im $i$-ten Glied ist schlicht $i$, von $i=2$ bis $i=12$:
\begin{equation}
  2+3+4+\dots+12 = \sum_{i=2}^{12} i .
\end{equation}
Ränder-Probe: $i=2$ gibt $2$ (erster Term), $i=12$ gibt $12$ (letzter
Term). \checkmark

\subsection*{b) $1+3+5+\dots+11$}
Ungerade Zahlen. Der $k$-te Summand ist $2k+1$. \emph{Startwert:} Der
erste Term ist $1$, also $2k+1=1\Rightarrow k=0$. \emph{Endwert:} Der
letzte Term $11$ entspricht $2k+1=11\Rightarrow k=5$:
\begin{equation}
  1+3+5+7+9+11=\sum_{k=0}^{5}(2k+1).
\end{equation}
\textbf{Achtung, der klassische Fehler: ein Start bei $k=1$} würde mit
$2\cdot1+1=3$ beginnen und den ersten Summanden $1$ verlieren. Die
Ränder-Probe fängt das.

\subsection*{c) $1+2^2+3^2+\dots+k^2$}
Der erste Term ist $1=1^2$, jeder Summand eine Quadratzahl $i^2$:
\begin{equation}
  1+2^2+3^2+\dots+k^2=\sum_{i=1}^{k} i^2 .
\end{equation}

\subsection*{d) $1+q^2+q^3+\dots+q^n$}
Neben der konstanten $1$ laufen die Terme $q^i$ von $i=2$ bis $i=n$
($q^1=q$ fehlt ja!):
\begin{equation}
  1+q^2+q^3+\dots+q^n = 1+\sum_{i=2}^{n} q^i .
\end{equation}
Ränder-Probe: der erste hinter dem $1$ ist $q^2$ (Index $2$), der letzte
$q^n$ (Index $n$). Ein Start bei $i=1$ wäre falsch, weil dann $q^1=q$
auftauchte, das in der Schreibweise nicht vorkommt.

\subsection*{e) $\sum_{i=1}^{12}(3i-1)$ ausgeschrieben}
Für $i=1,2,\dots,12$ ergeben sich $2,5,8,\dots,35$, also
\begin{equation}
  \sum_{i=1}^{12}(3i-1) = 2+5+8+\dots+35 .
\end{equation}
Erster Term $i=1$: $3-1=2$; letzter $i=12$: $36-1=35$. \checkmark

\subsection*{f) $\sum_{i=0}^{k}(i^3+1)$ ausgeschrieben}
\begin{equation}
  \sum_{i=0}^{k}(i^3+1)=1+2+9+28+\dots+(k^3+1) .
\end{equation}
Rechnung am Rand: $i=0$: $0+1=1$; $i=1$: $1+1=2$; $i=2$: $8+1=9$;
$i=3$: $27+1=28$. \checkmark

\begin{ergebnis}
\begin{align*}
  \text{a)}\ &\sum_{i=2}^{12} i, &
  \text{b)}\ &\sum_{k=0}^{5}(2k+1), &
  \text{c)}\ &\sum_{i=1}^{k} i^2,\\
  \text{d)}\ &1+\sum_{i=2}^{n}q^i, &
  \text{e)}\ &2+5+8+\dots+35, &
  \text{f)}\ &1+2+9+28+\dots+(k^3+1).
\end{align*}
\end{ergebnis}

\begin{kontrolle}
\textbf{Ränder-Probe in jedem Fall:}
\begin{itemize}
  \item b) $k=0$ gibt $2\cdot0+1=1$, $k=5$ gibt $11$ --- Anfang und Ende
        stimmen, also stimmt die ganze Summe.
  \item e) $i=1$ gibt $2$, $i=12$ gibt $35$.
  \item f) $i=0$ gibt $1$, $i=1$ gibt $2$, $i=2$ gibt $9$ --- die
        aufgeschriebenen Anfangswerte passen.
\end{itemize}

\textbf{Grenzfall (warum die Notation robust ist):} Ändert man nur einen
Rand, etwa Start bei $k=1$ in b), verschiebt sich die ganze Summe um ein
Glied --- ein einziger erster Term macht die gesamte Schreibweise falsch.
Die Ränder-Probe ist deshalb die kurze, sichere Abschlusspruefung.
\end{kontrolle}

\begin{merksatz}
Beim Umschreiben den Startwert so wählen, dass der erste Summand korrekt
geliefert wird; die Probe an \emph{beiden} Rändern (erster und letzter
Term) sichert die Schreibweise. Werte dürfen hier nicht ausgerechnet
werden, nur die Notation wechselt.
\end{merksatz}
